\section{Hard-Sibling Discriminative Attribute Selection}

The previous section described how CUB attributes are represented as semantic groups and used as auxiliary supervision in the proposed model. However, CUB-200-2011 contains 312 attributes \cite{wah2011cub}, and not all attributes are equally useful for distinguishing fine-grained bird species. Some attributes may be shared by many species, while others may be redundant or weakly discriminative. Therefore, this thesis studies an attribute selection strategy called \textit{Hard-Sibling Discriminative Attribute Selection} (HSDAS).

Although the original selection implementation can compute attribute scores at the species, family, and order levels, the proposed model in this thesis uses the species-level selected attributes for attribute-guided fusion. This choice is consistent with the architecture described in this chapter, where the attribute-guided module is attached to the species-level branch of H-CAST \cite{park2025visually}. Therefore, this section focuses only on species-level HSDAS.

\subsection{Motivation}

Fine-grained bird classification is challenging because visually similar species often belong to the same family. Such species may share global appearance, body shape, and dominant colors, but differ in localized visual traits such as bill shape, crown color, wing pattern, breast markings, or tail shape. A globally discriminative attribute may separate unrelated bird species, but may not help distinguish species that are taxonomically close. Therefore, the proposed selection strategy focuses on attributes that distinguish a species from its sibling species within the same family.

The key idea of HSDAS is to select attributes according to local taxonomic discrimination. For each species, the method compares its class-level attribute vector with the attribute vectors of sibling species. In addition, it gives special attention to hard siblings, which are sibling species that have similar attribute profiles. These hard siblings are likely to be difficult to distinguish, so attributes that separate them are expected to be useful for species-level recognition.

\subsection{Species-Level Attribute Representation}

Let the set of species be denoted as
\begin{equation}
    \mathcal{Y}^{s} = \{1, 2, \ldots, C_s\},
\end{equation}
where \(C_s = 200\) for CUB-200-2011. Each species \(c\) is associated with a class-level attribute vector
\begin{equation}
    \mathbf{a}_c =
    [a_{c,1}, a_{c,2}, \ldots, a_{c,312}]
    \in \mathbb{R}^{312}.
\end{equation}
In the implementation, the continuous CUB class-level attribute labels are normalized to the range \([0,1]\). Each element \(a_{c,j}\) indicates the class-level presence strength of attribute \(j\) for species \(c\), derived from the continuous CUB attribute annotations.

The taxonomy defines a mapping from species to family:
\begin{equation}
    \pi_f: \mathcal{Y}^{s} \rightarrow \mathcal{Y}^{f}.
\end{equation}
For a target species \(c\), its sibling species are defined as other species that share the same parent family:
\begin{equation}
    \mathcal{S}(c)
    =
    \{c' \in \mathcal{Y}^{s}
    \mid
    \pi_f(c') = \pi_f(c),\; c' \neq c \}.
\end{equation}

If a species has no sibling species, its attribute scores are set to zero and no attributes are selected from that species. In CUB, most species belong to families with multiple species, so sibling comparison is usually available.

\subsection{General Sibling Discriminativeness}

The first component of HSDAS measures how much stronger an attribute is for the target species compared with the average of its sibling species. For each species \(c\) and attribute \(j\), the mean sibling attribute value is computed as
\begin{equation}
    \bar{a}_{\mathcal{S}(c),j}
    =
    \frac{1}{|\mathcal{S}(c)|}
    \sum_{c' \in \mathcal{S}(c)}
    a_{c',j}.
\end{equation}

The general sibling discriminativeness score is then defined as
\begin{equation}
    d_{\text{all}}(c,j)
    =
    \max
    \left(
        a_{c,j} - \bar{a}_{\mathcal{S}(c),j},
        0
    \right).
\end{equation}

This score keeps only positive differences. If an attribute is not stronger for the target species than for its siblings on average, it is not considered discriminative for that species. This design makes the selected attributes species-specific: an attribute is useful only when it helps distinguish the target species from its local taxonomic neighborhood. Additionally, this formulation focuses on positive presence cues so attributes that are mainly discriminative through absence are not selected by this scoring rule.

\subsection{Hard-Sibling Discriminativeness}

The second component focuses on hard sibling species. Hard siblings are defined as sibling species whose attribute vectors are most similar to the target species. The similarity between species \(c\) and a sibling species \(c'\) is computed using cosine similarity:
\begin{equation}
    \mathrm{sim}(c,c')
    =
    \frac{
        \mathbf{a}_c^\top \mathbf{a}_{c'}
    }{
        \|\mathbf{a}_c\|_2 \|\mathbf{a}_{c'}\|_2
    }.
\end{equation}

For each species \(c\), the hard sibling set \(\mathcal{H}_r(c)\) contains the top \(r\) sibling species with the highest cosine similarity:
\begin{equation}
    \mathcal{H}_r(c)
    \subseteq
    \mathcal{S}(c).
\end{equation}
In the experiments of this thesis, the species-level setting uses \(r=3\).

For each hard sibling \(h \in \mathcal{H}_r(c)\), the positive attribute difference is computed as
\begin{equation}
    \Delta(c,h,j)
    =
    \max(a_{c,j} - a_{h,j}, 0).
\end{equation}

The hard-sibling discriminativeness score is a similarity-weighted sum of these positive differences:
\begin{equation}
    d_{\text{hard}}(c,j)
    =
    \sum_{h \in \mathcal{H}_r(c)}
    w_{c,h}
    \Delta(c,h,j),
\end{equation}
where the weight \(w_{c,h}\) is the normalized cosine similarity:
\begin{equation}
    w_{c,h}
    =
    \frac{
        \mathrm{sim}(c,h)
    }{
        \sum_{h' \in \mathcal{H}_r(c)}
        \mathrm{sim}(c,h') + \epsilon
    }.
\end{equation}

This weighting gives more influence to sibling species that are more similar to the target species. As a result, the method emphasizes attributes that separate a species from its most confusing relatives.

\subsection{Final Attribute Score}

The final HSDAS score combines the general sibling discriminativeness and the hard-sibling discriminativeness:
\begin{equation}
    s(c,j)
    =
    d_{\text{all}}(c,j)
    +
    d_{\text{hard}}(c,j).
\end{equation}

The first term captures attributes that distinguish the target species from all sibling species on average. The second term captures attributes that distinguish the target species from the most similar sibling species. Combining both terms allows the method to balance general local discrimination and hard-case discrimination.

For each species, attributes with zero final score are removed from consideration. The remaining attributes are sorted in descending order by three criteria:
\begin{equation}
    s(c,j), \quad d_{\text{hard}}(c,j), \quad d_{\text{all}}(c,j).
\end{equation}
The final score is the primary ranking criterion. The hard-sibling score is used as the first tie-breaker, giving priority to attributes that help with difficult sibling distinctions. The general sibling score is used as the second tie-breaker.

\subsection{Top-\(k\) Selection and Attribute Union}

After scoring all attributes for each species, HSDAS selects the top \(k\) attributes for each species. In this thesis, the species-level setting uses \(k=5\). Let \(\mathcal{T}_k(c)\) denote the top-\(k\) selected attributes for species \(c\):
\begin{equation}
   \mathcal{T}_k(c)
    =
    \operatorname{TopK}_{k}
    \left(
    \{j \mid s(c,j) > 0\},\; s(c,j)
    \right).
\end{equation}

The final selected attribute set is obtained by taking the union of the selected attributes over all species:
\begin{equation}
    \mathcal{M}
    =
    \bigcup_{c \in \mathcal{Y}^{s}}
    \mathcal{T}_k(c).
\end{equation}

This union operation produces a compact global attribute subset for the dataset. Although attributes are selected based on species-level sibling discrimination, the final mask is shared by the attribute-guided model. In the main setting of this thesis, this process produces the HSDAS-199 attribute subset.

The selected attributes are then converted into a group mask according to the 28 CUB attribute groups. For each group \(g\), the mask contains only the selected attribute indices that belong to that group:
\begin{equation}
    M_g
    =
    \mathcal{M} \cap I_g,
\end{equation}
where \(I_g\) is the full index set of attributes in group \(g\). These group masks define the active attribute states used by the group-specific attribute classifiers.

\subsection{Algorithm Summary}

Algorithm~\ref{alg:hsdas_species} summarizes the species-level HSDAS procedure used in this thesis. The detailed definitions of \(d_{\mathrm{all}}\), \(d_{\mathrm{hard}}\), and \(s(c,j)\) are given in the previous subsections.

\begin{algorithm}[t]
\footnotesize
\caption{Species-Level Hard-Sibling Discriminative Attribute Selection}
\label{alg:hsdas_species}
\begin{algorithmic}[1]
\REQUIRE Attribute matrix \(A \in \mathbb{R}^{C_s \times 312}\), species-to-family mapping \(\pi_f\), number of hard siblings \(r\), number of selected attributes per species \(k\)
\ENSURE Selected attribute set \(\mathcal{M}\)

\STATE Initialize \(\mathcal{M} \leftarrow \emptyset\)

\FOR{each species \(c \in \mathcal{Y}^{s}\)}
    \STATE Find sibling species:
    \(\mathcal{S}(c) \leftarrow \{c' \in \mathcal{Y}^{s} \mid \pi_f(c')=\pi_f(c), c' \neq c\}\)

    \IF{\(\mathcal{S}(c)=\emptyset\)}
        \STATE Continue
    \ENDIF

    \STATE Compute mean sibling attribute vector:
    \(\bar{\mathbf{a}}_{\mathcal{S}(c)} \leftarrow
    |\mathcal{S}(c)|^{-1}\sum_{c' \in \mathcal{S}(c)} \mathbf{a}_{c'}\)

    \STATE Compute \(d_{\mathrm{all}}(c,j)\) for all attributes \(j=1,\ldots,312\)

    \STATE Compute cosine similarity between \(c\) and each sibling \(c' \in \mathcal{S}(c)\)

    \STATE Select hard siblings \(\mathcal{H}_r(c)\) as the top-\(r\) most similar species in \(\mathcal{S}(c)\)

    \STATE Compute normalized hard-sibling weights \(w_{c,h}\) for all \(h \in \mathcal{H}_r(c)\)

    \STATE Compute \(d_{\mathrm{hard}}(c,j)\) and final score
    \(s(c,j)=d_{\mathrm{all}}(c,j)+d_{\mathrm{hard}}(c,j)\) for all \(j\)

    \STATE Select top-\(k\) positive-score attributes:
    \[
        \mathcal{T}_k(c)
        \leftarrow
        \operatorname{TopK}_{k}
        \left(
            \{j \mid s(c,j)>0\},
            s(c,j)
        \right)
    \]

    \STATE Update selected set:
    \(\mathcal{M} \leftarrow \mathcal{M} \cup \mathcal{T}_k(c)\)
\ENDFOR

\RETURN \(\mathcal{M}\)
\end{algorithmic}
\end{algorithm}

\subsection{Use in the Proposed Model}

The selected HSDAS attribute set is used to build the attribute mask for the proposed attribute-guided H-CAST model. This mask determines which attribute states are retained in each attribute group and therefore controls the output dimensions of the group-specific attribute classifiers.

Importantly, HSDAS is used only to construct the training supervision mask. It does not introduce additional inputs during inference. At test time, the model still receives only an image. The selected attributes influence the model indirectly by supervising the attribute-guided module during training and by shaping the attribute-oriented features used for species-level fusion.

This design has two advantages. First, it reduces the attribute space from the full 312 attributes to a compact subset focused on species-level discrimination. Second, it aligns attribute selection with the fine-grained nature of the task by prioritizing attributes that distinguish sibling and hard-sibling species. Therefore, HSDAS provides a taxonomy-aware attribute subset that is more directly connected to the objective of species-level bird classification than a purely global attribute selection criterion.